\chapter{Myoclonus}

\section*{Definition}
Myoclonus refers to sudden, brief, involuntary, shock-like jerks of a single muscle or group of muscles, caused by abrupt muscular contraction (positive myoclonus) or sudden cessation of contraction (negative myoclonus or asterixis). It may be focal, multifocal, generalized, or segmental and is classified by causes, distribution, and neurophysiologic origin.

\section*{What Happens in Your Body}
Myoclonus arises from hyperexcitable neuronal populations at various levels of the neuraxis. Cortical myoclonus involves hyperexcitable sensorimotor cortex with abnormally enhanced long-loop reflexes and giant somatosensory evoked potentials. Subcortical myoclonus originates from brainstem reticular formation, particularly the caudal medulla (reticular reflex myoclonus). Spinal myoclonus is generated by hyperexcitable spinal interneurons after segmental deafferentation or injury. Negative myoclonus (asterixis) results from sudden loss of muscle tone due to impaired sensorimotor integration, commonly in metabolic encephalopathy.

\section*{Causes}
\begin{itemize}
  \item Physiologic myoclonus: Sleep starts (hypnic jerks), hiccups, exercise-induced myoclonus, anxiety-related jerks
  \item Epileptic myoclonus: Juvenile myoclonic epilepsy, progressive myoclonic epilepsies (Unverricht-Lundborg, Lafora disease, MERRF, sialidosis), myoclonic-astatic epilepsy (Doose syndrome)
  \item Symptomatic (secondary) myoclonus: Post-hypoxic (Lance-Adams syndrome), metabolic encephalopathy (hepatic, renal, respiratory failure), electrolyte disturbance (hyperosmolar state, hypocalcemia), drug toxicity (lithium, opioids, tricyclic antidepressants, levodopa), toxin exposure (heavy metals, organophosphates)
  \item Neurodegenerative causes: Creutzfeldt-Jakob disease, progressive supranuclear palsy, multiple system atrophy, corticobasal degeneration, Alzheimer disease, Huntington disease, dentatorubral-pallidoluysian atrophy
  \item Infectious and paraneoplastic: Viral encephalitis (herpes simplex, HIV, arboviruses), subacute sclerosing panencephalitis, paraneoplastic syndromes (anti-CV2, anti-amphiphysin, anti-Ri)
  \item Storage diseases: Gaucher disease (type 3), Tay-Sachs disease, Niemann-Pick type C, ceroid lipofuscinosis
\end{itemize}

\section*{When to See a Doctor}
See your doctor if:
\begin{itemize}
  \item New or progressive myoclonus, especially with cognitive decline, seizures, or gait impairment
  \item Myoclonus following cardiac arrest or hypoxic event (Lance-Adams syndrome)
  \item Multifocal or generalized myoclonus with altered mental status (possible Creutzfeldt-Jakob disease or metabolic encephalopathy)
  \item Myoclonus in the setting of known or suspected malignancy
  \item Family history of myoclonus, epilepsy, or neurodegenerative disease
\end{itemize}

\section*{Related Symptoms}
\begin{itemize}
  \item Seizures
  \item Ataxia
  \item Tremor
  \item Chorea
  \item Dystonia
  \item Cognitive decline
  \item Dysarthria
  \item Asterixis
\end{itemize}
\textbf{Important:} Action myoclonus (provoked by voluntary movement) is particularly disabling and characteristic of post-hypoxic myoclonus and progressive myoclonic epilepsies; electrophysiology (EEG-EMG back-averaging) localizes the myoclonus generator.

\section*{Self-Care / Home Management}
\begin{itemize}
  \item Identify and avoid triggers such as sudden movements, bright lights, or startle stimuli.
  \item Ensure safety during jerks by avoiding activities requiring fine motor control during exacerbations.
  \item Use padded protective gear (helmets, joint guards) if falls occur due to negative myoclonus.
  \item Maintain consistent sleep schedules as sleep deprivation worsens myoclonus.
\end{itemize}

\section*{How Your Doctor Will Evaluate You}
\begin{itemize}
  \item Clinical characterization: Body distribution, stimulus sensitivity, relation to action or posture, and temporal pattern.
  \item Electrophysiologic studies: EEG-EMG polygraphy, jerk-locked back-averaging, and somatosensory evoked potentials identify cortical, subcortical, or spinal origin.
  \item Laboratory evaluation: Comprehensive metabolic panel, ammonia, liver and renal function, toxicology screen, thyroid function, and infectious serologies.
  \item Neuroimaging (brain MRI) and CSF analysis (14-3-3 protein, neuron-specific enolase for suspected prion disease) as indicated.
\end{itemize}

\section*{Treatment Options}
\begin{itemize}
  \item Correct underlying metabolic derangements (dialysis for uremia, lactulose for hepatic encephalopathy, electrolyte repletion).
  \item Remove or adjust offending medications in drug-induced myoclonus.
  \item First-line pharmacotherapy: Levetiracetam, valproic acid, and clonazepam are most effective across myoclonus subtypes.
  \item Second-line agents: Piracetam (particularly for post-hypoxic myoclonus), topiramate, zonisamide, primidone, and 5-hydroxytryptophan.
  \item Deep brain stimulation (pallidal or thalamic) may benefit medically refractory cases in selected degenerative disorders.
\end{itemize}

\section*{References}
\begin{thebibliography}{99}
\bibitem{myoclonus:ref1} Caviness JN. ``Myoclonus.'' \textit{Mayo Clinic Proceedings}, 2009;84(9):833--846.
\bibitem{myoclonus:ref2} Shibasaki H, Hallett M. ``Electrophysiological studies of myoclonus.'' \textit{Muscle \& Nerve}, 2005;31(2):157--174.
\bibitem{myoclonus:ref3} Marsden CD, Hallett M, Fahn S. ``The nosology and pathophysiology of myoclonus.'' \textit{Movement Disorders}, 2012;27(4):479--492.
\bibitem{myoclonus:ref4} Levy A, Chen R. ``Myoclonus: pathophysiology and treatment options.'' \textit{Current Treatment Options in Neurology}, 2016;18(5):21.
\bibitem{myoclonus:ref5} Caviness JN, Brown P. ``Myoclonus: current concepts and recent advances.'' \textit{Lancet Neurology}, 2004;3(10):598--607.
\bibitem{myoclonus:ref6} Obeso JA, Rothwell JC, Marsden CD. ``The spectrum of cortical myoclonus.'' \textit{Brain}, 1985;108(1):193--224.

\end{thebibliography}
